\documentclass{article}
\usepackage{amsmath,amssymb,amsthm}
\usepackage{algorithm}
\usepackage{algorithmic}
\usepackage{hyperref}
\usepackage{enumitem}
\newtheorem{assumption}{Assumption}
\newtheorem{theorem}{Theorem}
\newtheorem{lemma}{Lemma}
\newcommand{\argmin}{\mathop{\mathrm{arg\,min}}}
\newcommand{\grad}{\nabla}
\newcommand{\Breg}{D_{h}}
\newcommand{\prox}{\operatorname{prox}}
\newcommand{\norm}[1]{\left\|#1\right\|}
\begin{document}

\section*{Algorithm Name}
\textbf{Relative (Mirror) Gradient Descent (RGD)}  

The method is a Bregman‑proximal gradient step where the objective $f$ is assumed to be \emph{smooth relative to} a strongly convex distance‑generating function $h$.

\section*{Mathematical Formulation}
Let $h:\mathbb R^{d}\to\mathbb R$ be differentiable and $\mu$‑strongly convex w.r.t. a norm $\|\cdot\|$, i.e.  
\[
h(y)\ge h(x)+\langle \grad h(x),y-x\rangle+\frac{\mu}{2}\|y-x\|^{2},
\qquad\forall x,y.
\]
Its Bregman divergence is  
\[
\Breg (y,x)=h(y)-h(x)-\langle \grad h(x),y-x\rangle .
\]

Given a stepsize $\eta_k>0$, the iteration reads
\begin{equation}
\boxed{
x_{k+1}= \argmin_{x\in\mathbb R^{d}}
\Big\{ \langle \grad f(x_k),\,x\rangle+\frac{1}{\eta_k}\,\Breg (x,x_k) \Big\}
}\tag{1}
\end{equation}
The optimality condition of (1) is
\begin{equation}
\grad h(x_{k+1}) = \grad h(x_k)-\eta_k \grad f(x_k). \tag{2}
\end{equation}
Define the \emph{gradient mapping} (relative to $h$) as
\begin{equation}
G_{\eta_k}(x_k):=\frac{1}{\eta_k}\bigl(x_k-\prox_{\eta_k h}\bigl(x_k-\eta_k\grad f(x_k)\bigr)\bigr)
               =\frac{1}{\eta_k}\bigl(x_k-x_{k+1}\bigr). \tag{3}
\end{equation}

\section*{Pseudocode}
\begin{algorithm}[H]
\caption{Relative Gradient Descent (RGD)}\label{alg:RGD}
\begin{algorithmic}[1]
\STATE \textbf{Input:} initial point $x_0\in\mathbb R^{d}$, stepsizes $\{\eta_k\}_{k\ge0}>0$, prox‑function $h$.
\FOR{$k=0,1,2,\dots$}
\STATE Compute stochastic (here deterministic) gradient $g_k=\grad f(x_k)$.
\STATE $x_{k+1}\gets\argmin_{x}\Big\{\langle g_k,x\rangle+\frac{1}{\eta_k}\Breg(x,x_k)\Big\}$.
\IF{termination criterion satisfied (e.g. $\|G_{\eta_k}(x_k)\|\le\varepsilon$)} 
\STATE \textbf{break}
\ENDIF
\ENDFOR
\STATE \textbf{Output:} $x_{K}$ (or the best iterate).
\end{algorithmic}
\end{algorithm}

\section*{Dry Run Example}
Consider the scalar problem  
\[
f(w)=(w-3)^2,\qquad h(w)=\tfrac12 w^{2}\;\;\Longrightarrow\;\;
\Breg (w,w_k)=\tfrac12 (w-w_k)^2 .
\]
We take $w_0=0$ and a constant stepsize $\eta=0.1$.

\bigskip
\noindent\textbf{Iteration $k=0$}
\[
\grad f(w_0)=2(w_0-3)=-6.
\]
Optimality condition (2) gives  
\[
\grad h(w_1)=\grad h(w_0)-\eta\grad f(w_0)
\;\Longrightarrow\; w_1 = w_0-2\eta (w_0-3)
                =0-2\cdot0.1\cdot(-3)=0.6 .
\]
Objective value $f(w_1)=(0.6-3)^2=5.76$.

\bigskip
\noindent\textbf{Iteration $k=1$}
\[
\grad f(w_1)=2(0.6-3)=-4.8.
\]
\[
w_2 = w_1-2\eta (w_1-3)
    =0.6-2\cdot0.1\cdot(-2.4)=0.6+0.48=1.08 .
\]
$f(w_2)=(1.08-3)^2=3.6864$.

\bigskip
\noindent\textbf{Iteration $k=2$}
\[
\grad f(w_2)=2(1.08-3)=-3.84.
\]
\[
w_3 = w_2-2\eta (w_2-3)
    =1.08-2\cdot0.1\cdot(-1.92)=1.08+0.384=1.464 .
\]
$f(w_3)=(1.464-3)^2=2.3660$.

The iterates move monotonically toward the minimiser $w^\star=3$.

\newpage
\section*{Assumptions}
\begin{assumption}[Relative Smoothness]\label{ass:smooth}
There exists $L>0$ such that for all $x,y\in\mathbb R^{d}$
\[
f(y)\le f(x)+\langle \grad f(x),y-x\rangle+L\,\Breg (y,x).
\tag{A1}
\]
\end{assumption}
\begin{assumption}[Strong Convexity of the Prox‑function]\label{ass:strong}
The distance‑generating function $h$ is $\mu$‑strongly convex ($\mu>0$) w.r.t. $\|\cdot\|$, i.e.
\[
\Breg (y,x)\ge\frac{\mu}{2}\|y-x\|^{2},\qquad\forall x,y.
\tag{A2}
\]
\end{assumption}
\begin{assumption}[Bounded Below]\label{ass:bound}
The objective $f$ is bounded below: $f^\star:=\inf_{x}f(x)>-\infty$.
\tag{A3}
\end{assumption}
\begin{assumption}[Stepsize]\label{ass:stepsize}
The stepsizes satisfy $0<\eta_k\le \frac{1}{L}$ for all $k$.
\tag{A4}
\end{assumption}

\section*{Main Convergence Theorem}
\begin{theorem}[Convergence of RGD for Non‑Convex $f$]\label{thm:main}
Under Assumptions~\ref{ass:smooth}–\ref{ass:stepsize}, let $\{x_k\}$ be generated by \eqref{alg:RGD}. Define the gradient mapping \eqref{3}. Then for any integer $K\ge1$,
\[
\frac{1}{K}\sum_{k=0}^{K-1}\norm{G_{\eta_k}(x_k)}^{2}
\;\le\;
\frac{2L}{\mu}\,\frac{f(x_0)-f^\star}{K}.
\tag{4}
\]
Consequently,
\[
\min_{0\le k\le K-1}\norm{G_{\eta_k}(x_k)}^{2}
\;\le\;
\frac{2L}{\mu}\,\frac{f(x_0)-f^\star}{K}.
\]
\end{theorem}

\section*{Theoretical Properties}
\begin{itemize}[leftmargin=*]
\item \textbf{Descent guarantee.} Inequality~\eqref{4} shows that the average squared norm of the gradient mapping decays at the rate $O(1/K)$, which is the optimal rate for first‑order methods on smooth non‑convex problems.
\item \textbf{Stepsize sensitivity.} The condition $\eta_k\le 1/L$ is necessary to keep the descent term positive; any larger $\eta_k$ may break the relative smoothness bound.
\item \textbf{Relation to Mirror Descent.} When $h(x)=\tfrac12\|x\|^{2}$, $\Breg$ reduces to the Euclidean squared distance and $G_{\eta_k}$ becomes the usual gradient step $G_{\eta_k}(x_k)=\grad f(x_k)$. Hence the theorem recovers the classical $O(1/K)$ guarantee for (non‑convex) gradient descent.
\item \textbf{Comparison to baselines.} Compared with Euclidean SGD (Example~1) the bound contains the factor $L/\mu$, reflecting the geometry induced by $h$. For ill‑conditioned problems a judicious choice of $h$ (e.g., a barrier for the feasible set) can make $\mu$ large and improve the constant.
\end{itemize}

\section*{Discussion}
The result hinges on relative smoothness rather than ordinary Lipschitz continuity of $\grad f$. This permits handling objectives such as logistic loss with the Kullback–Leibler Bregman distance. The theorem is deterministic; stochastic extensions follow by taking expectations and adding a variance term, analogous to the SGD proof.

\newpage
\section*{Preliminaries (Definitions and Lemmas)}
\begin{lemma}[Three‑Point Identity for Bregman Divergence]\label{lem:3pt}
For any $x,y,z$,
\[
\Breg (x,z)-\Breg (y,z)-\Breg (x,y)=\langle \grad h(x)-\grad h(y),\,z-y\rangle .
\tag{5}
\]
\end{lemma}
\begin{proof}
By definition,
\[
\Breg (x,z)=h(x)-h(z)-\langle \grad h(z),x-z\rangle,
\]
\[
\Breg (y,z)=h(y)-h(z)-\langle \grad h(z),y-z\rangle,
\]
\[
\Breg (x,y)=h(x)-h(y)-\langle \grad h(y),x-y\rangle .
\]
Subtracting the second and third from the first and rearranging yields \eqref{5}. ∎
\end{proof}

\section*{Main Proof (Step‑by‑Step Derivation)}
We prove Theorem~\ref{thm:main}. All steps are numbered for easy reference.

\bigskip
\noindent\textbf{[Step 1]} \emph{Relative smoothness inequality.}  
From Assumption~\ref{ass:smooth} with $x=x_k$, $y=x_{k+1}$,
\[
f(x_{k+1})\le f(x_k)+\langle \grad f(x_k),x_{k+1}-x_k\rangle+L\,\Breg (x_{k+1},x_k).
\tag{6}
\]

\noindent\textbf{[Step 2]} \emph{Optimality condition of the subproblem.}  
The definition \eqref{1} implies that $x_{k+1}$ minimizes a convex function; the first‑order optimality condition is exactly \eqref{2}:
\[
\grad h(x_{k+1}) = \grad h(x_k)-\eta_k\grad f(x_k).
\tag{7}
\]

\noindent\textbf{[Step 3]} \emph{Express the linear term in \eqref{6} via Bregman divergences.}  
Take the inner product of \eqref{7} with $x_{k+1}-x_k$:
\[
\langle \grad h(x_{k+1})-\grad h(x_k),x_{k+1}-x_k\rangle
   = -\eta_k \langle \grad f(x_k),x_{k+1}-x_k\rangle .
\tag{8}
\]
Rearrange \eqref{8} to obtain
\[
\langle \grad f(x_k),x_{k+1}-x_k\rangle
   = -\frac{1}{\eta_k}\langle \grad h(x_{k+1})-\grad h(x_k),x_{k+1}-x_k\rangle .
\tag{9}
\]

\noindent\textbf{[Step 4]} \emph{Apply the three‑point identity \eqref{5}.}  
Set $(x,y,z)=(x_{k+1},x_k,x_k)$ in \eqref{5}:
\[
\Breg (x_{k+1},x_k)-\Breg (x_k,x_k)-\Breg (x_{k+1},x_k)
   =\langle \grad h(x_{k+1})-\grad h(x_k),x_k-x_k\rangle=0 .
\]
A more useful choice is $(x,y,z)=(x_{k+1},x_k,x_k)$ which yields
\[
\langle \grad h(x_{k+1})-\grad h(x_k),x_{k+1}-x_k\rangle
   =\Breg (x_{k+1},x_k)+\Breg (x_k,x_{k+1}) .
\tag{10}
\]
Because $\Breg (x_k,x_{k+1})\ge \frac{\mu}{2}\|x_{k+1}-x_k\|^{2}$ by strong convexity (Assumption~\ref{ass:strong}), we have
\[
\langle \grad h(x_{k+1})-\grad h(x_k),x_{k+1}-x_k\rangle
   \ge \Breg (x_{k+1},x_k)+\frac{\mu}{2}\|x_{k+1}-x_k\|^{2}.
\tag{11}
\]

\noindent\textbf{[Step 5]} \emph{Plug \eqref{9}–\eqref{11} into \eqref{6}.}  
From \eqref{9} and \eqref{11},
\[
\langle \grad f(x_k),x_{k+1}-x_k\rangle
   \le -\frac{1}{\eta_k}\Bigl(\Breg (x_{k+1},x_k)+\frac{\mu}{2}\|x_{k+1}-x_k\|^{2}\Bigr).
\tag{12}
\]
Insert \eqref{12} into \eqref{6}:
\[
f(x_{k+1})\le f(x_k)
   -\frac{1}{\eta_k}\Breg (x_{k+1},x_k)
   -\frac{\mu}{2\eta_k}\|x_{k+1}-x_k\|^{2}
   +L\,\Breg (x_{k+1},x_k).
\tag{13}
\]

\noindent\textbf{[Step 6]} \emph{Collect the Bregman terms.}  
Since $\eta_k\le 1/L$ (Assumption~\ref{ass:stepsize}), we have $L-\frac{1}{\eta_k}\le 0$. Hence
\[
-\frac{1}{\eta_k}\Breg (x_{k+1},x_k)+L\,\Breg (x_{k+1},x_k)
   = -\Bigl(\frac{1}{\eta_k}-L\Bigr)\Breg (x_{k+1},x_k)
   \le 0 .
\tag{14}
\]
Discarding this non‑positive term from \eqref{13} yields a clean descent inequality:
\[
f(x_{k+1})\le f(x_k)-\frac{\mu}{2\eta_k}\|x_{k+1}-x_k\|^{2}.
\tag{15}
\]

\noindent\textbf{[Step 7]} \emph{Express the distance via the gradient mapping.}  
From definition \eqref{3},
\[
\|G_{\eta_k}(x_k)\|^{2}
   =\frac{1}{\eta_k^{2}}\|x_k-x_{k+1}\|^{2}
   =\frac{1}{\eta_k^{2}}\|x_{k+1}-x_k\|^{2}.
\tag{16}
\]
Rearranging \eqref{15} and using \eqref{16} gives
\[
\|G_{\eta_k}(x_k)\|^{2}
   \le \frac{2L}{\mu}\,\frac{f(x_k)-f(x_{k+1})}{\eta_k}.
\tag{17}
\]

\noindent\textbf{[Step 8]} \emph{Sum over $k=0,\dots,K-1$.}  
Assuming a constant stepsize $\eta_k\equiv\eta\le 1/L$, sum \eqref{17}:
\[
\sum_{k=0}^{K-1}\|G_{\eta}(x_k)\|^{2}
   \le \frac{2L}{\mu\eta}\sum_{k=0}^{K-1}\bigl(f(x_k)-f(x_{k+1})\bigr)
   =\frac{2L}{\mu\eta}\bigl(f(x_0)-f(x_K)\bigr)
   \le\frac{2L}{\mu\eta}\bigl(f(x_0)-f^\star\bigr).
\tag{18}
\]

\noindent\textbf{[Step 9]} \emph{Divide by $K$ and replace $\eta$ by its upper bound $1/L$.}  
Because $\eta\le 1/L$, we have $1/(\eta)\le L$, thus
\[
\frac{1}{K}\sum_{k=0}^{K-1}\|G_{\eta}(x_k)\|^{2}
   \le \frac{2L}{\mu}\,\frac{f(x_0)-f^\star}{K}.
\tag{19}
\]
Inequality \eqref{19} is exactly the bound \eqref{4} stated in Theorem~\ref{thm:main}. The claim about the minimum of the sequence follows from the fact that the average upper‑bounds the minimum.

\noindent\textbf{[Step 10]} \emph{Conclusion.}  
Thus the algorithm generates a sequence whose gradient‑mapping norm decays at the rate $O(1/K)$, completing the proof. ∎

\section*{Rate Derivation (With Constants)}
The constant in \eqref{4} is $\displaystyle C=\frac{2L}{\mu}\bigl(f(x_0)-f^\star\bigr)$.  
If $L$ and $\mu$ are known, one may select the maximal admissible stepsize $\eta=1/L$ to minimise $C$.

\section*{Special Cases}
\begin{enumerate}
\item \textbf{Euclidean setting.} $h(x)=\tfrac12\|x\|^{2}$ gives $\mu=1$, $\Breg (y,x)=\tfrac12\|y-x\|^{2}$, and $G_{\eta}(x)=\grad f(x)$. Theorem~\ref{thm:main} reduces to the classic $O(1/K)$ bound for gradient descent on $L$‑smooth non‑convex functions.
\item \textbf{KL‑divergence.} For $h(x)=\sum_{i}x_i\log x_i -x_i$, $\mu=1$ on the simplex and $\Breg$ becomes the Kullback–Leibler divergence. The theorem then applies to problems such as Poisson regression where ordinary Lipschitz smoothness fails but relative smoothness holds.
\item \textbf{Adaptive stepsizes.} If one uses $\eta_k=1/(L_k)$ with a locally estimated $L_k$, the same proof works provided $\eta_k\le1/L$ for each $k$.
\end{enumerate}

\end{document}